\documentclass[]{../../../../tex/classes/styledArticle}
\usepackage{../../../../tex/packages/hebrewSupport}

\author{שחר פרץ}
\title{\textit{היסטוריה 44}}
\begin{document}
	\maketitle
	\section*{גורל יהודי ארצות האסלם}
	
	המספר של המקרן: 28757משהו
	
	מחר יש חרא צהלי, שבוע הבא מבחן בלשון, ועוד שבועיים יום נגישות. 
	
	ארצות המזרח – מצרים לוב סוריה לבנון עיראק תוניס אלז'יר וכו'. יש מדינות מוסלמיות לא רק במזרח התיכון וצפון אפריקה (כמו צרפת לדוגמה) אבל אנחנו לא מדברים עליהם. 
	
	נרצה לדבר כמעה על מעמד היהודים לפני הגעת המעצמות הקולוניאליות. מעמדם אז היה כמעמדם של היהודים באסלם. ליהודים ולנוצרים יש מעמד באסלם, הקרוי ''עמי הספר`` בתרגום. האסלם דת מונוטאיסטית מורחבת יותר מהיהדות והנצרות, והיא מתבססת על שתי הדתות בהרבה אלמנטים. המשמעות של אותו המעמד, שלא אמורים להרוג את היהודים. הוא נחות יותר ביחס לאסלמים. לדוגמה, בית כנסת לא יכול להיות גבוה, אסור ליהודים לרכב על סוס או חמור (כי אז הם נהיים גבוהים יותר). במהלך השנים נוצרו יותר ויותר חוקים (לדוגמה, היהודים שילמו מס כדי שיגינו עליהם). חשוב להבין שהמעמד של היהודים באסלם יותר טוב מזה של הנוצרים, ואפילו יותר טוב מזה של הנצרות. לאורך ההיסטוריה, מצבם של היהודים במדינות מוסלמיות היה הרבה יותר טוב מזה של היהודים במדינות נוצריות. העובדה שאנו חיים בעולם שבו במאת השנים האחרונות האסלם הוא זה שפוגע ביהודים, לא אומר שכך זה היה. הגירושים, האינקווזציה והפוגרומים, נעשו ע''י הנוצרים. יתרה מכך, הדתיים הנוצריים האדוקים רואים באירועים האילו פגיעה באסלם. 
	
	היהודים באותן המדינות חיו בדרך־כלל בשכונות נפרדות, בדומה לגטאות באירופה. הם שמרו על אורך חייהם. בשלב מסוים, השינויים שקרו באירופה הגיעו למדינות המוסלמיות, לדוגמה גידול דמוגרפי, בריאות, יציאה מהגטאות, וכו'. ב־1948 חיו בין 800K למיליון יהודים (''תלוי מי סופר ובאיזו נקודת זמן``). ב־67 היו 200K. 
	
	בעמוד כלשהו בספר יש לנו טבלה שמראה את ההבדלים. באיראן מצב של היהודים היה טוב – השאה היה טוב ליהודים. איראן התנהלה בצורה אחרת, יותר מערבית וטכנולוגית. שלטון השאה אפשר להיהודים להגיע לשכבות הגבוהות מבחינה חברתית. אנשים מאיראן יוכלו לספר כמה הם היו עשירים לפני המהפכה, וכמה כסף הם השאירו שם. ואכן, 80 אלף היהודים באיראן ב־48 נשארו שם עד 69. 
	
	הערת צד: יש יהודים באיראן. לא הרבה, אבל יש. מצבם יחסית סביל. מבחינה דתית ליהודים באיראן יש אפשרות לנהל את חיי דתם. בפסח יש מצות, הן פשוט לא מגיעות ממדינות עוינות. איראן ברמה העקרונית היא אנטי־ציונית/ישראלית ולא אנטי־יהודית. 
	
	בכל מקרה רוב הקהילות היהודיות נעלמות ומחוסלות. למעלה מ־80\% מהיהודים הגרו, בעיקר לארץ. יש לציין בעיקר את הסיבות הבאות: 
	\begin{itemize}
		\item \textit{הדה־קולוניזציה: }שתי סיבות: 
		\begin{itemize}
			\item שינוי תרבותי: בתפיסה שלהם, היהודים היו יותר קרובים לנוצרים. לכן היהודים הפכו להיות חלק מהשלטון הרבה פעמים. לכן נוצר מצב שיש שינוי תרבותי – היהודים יותר קרובים לשלטון מאשר לאוכלוסיה המקומית. 
			\item התגברות השנאה אליהם: כאשר המקומיים נלחמים בשלטון הכובש כדי לסלקו, השנאה ליהודים שזוהו עם השלטון גברה. לכן היו הרבה יותר פגיעות ביהודים בתקופת הדה־קולוניזציה. 
		\end{itemize}
		\item \textit{הקמת מדינת ישראל ומלחמת העצמאות: }רוב העליה ממדינות האסלם, מתרחשת. מדינת ישראל עושה מהלך שנקרא חיסול גלויות. בהבט הזה נעשה הבחנה בין שני מעגלים של מדינות: 
		\begin{itemize}
			\item מעגל ראשון – מדינות שנלחמות ישירות במדינת ישראל. לבנון וירדן לא רלוונטיות כי אין שם יהודים, אבל בעיראק, סוריה ובמצרים יש יהודים רבים. במדינות אלו יש פגיעה מאוד רעה ביהודים, והם נמצאים לרוב בסכנת חיים מידית וממשית. 
			\item המעגל השני – יש השפעה על היהודים, בעוצמה פחותה. יוצאת הדופן הייתה לוב בה ביוני 1948 היה גל אלימות קשה. 
		\end{itemize}
		\item \textit{הפאן ערביות והליגה הערבית: }פאן ערביות = שאיפה ללאום ערבי אחד. וותאניה משמעה לאומיות מקומית (לדוגמה לאומיות פלסטינית, עיראקית וכו'). קראמייה היא הפאן ערביות, העולם הערבי. התפיסה שהעולם הערבי כמכלול, הוא ביחד. כאשר התפיסה התגברה, נוצרה הליגה הערבית, ונוצרה התחושה של העול םהערבי כולו מול מדינת ישראל. מכאן שזה לא משנה אם אנו נלחמים באופן ישיר או לא, העולם הערבי כולו נגד ישראל וצריך לסייע לפלסטינים. ככל שהתפיסה התפתחה, הפגיעה ביהודים במדינות מעגל שני הלכה וגדלה. 
		\item \textit{ציונות וזיקה דתית לארץ: }הארגונים הציונים צוברים תאוצה משנות ה־30 ובעיקר אחרי מלחמת העולם. הזיקה והרצון להגיע לא''י מסיבות של ציונות התגברה. היו גם כאלו שהזיקה שלהם הייתה זיקה דתית, אם כי הם היו מיעוט. 
	\end{itemize}
	
	חשוב להפריד בין מוסלמים לערבים. ערבים זה מוצא אתני, מוסלמים זו דת. זה לא אותו הדבר. לדוגמה טוריקה מדינה מוסלמית שאיננה ערבית. הפאן ערביות לא מאחדת של המוסלמים לתוך הסיפור הזה. 
	
	תפישת מדינת ישראל היא כמדינת היהודים, בארץ ובגולה. היא המקור, הפתרון, והאחריות. לדוגמה, כאשר יש אנטישמיות בעולם, מדינת ישראל צריכה לסייע ליהודים בעולם. יהודים שנקלעו למלחמה באוקראינה או לאנטישמיות בצרפת יכלו לעלות לארץ. 
	
	אז איך עלו כל היהודים האלו לארץ? מבצעים של המוסד. המוסד מארגן מטוסים, תשתיות להעביר את היהודים וכו'. היו מדינות שבהן היה אפשר לבצע את כל הסיפור מול השלטון, ובמקרים אחרים העבירו אותם למדינות אחרות והעלו אותם משם. המטרה העיקרית הייתה להציל את הקהילות היהודיות נמצאות בסכנה. היו מדינות בהן עשו עליה סלקטיבית (מי שלא יכול לעבוד ולפרנס, ישאר באופן זמני במדינות האם, בהינתן שהם אינם נמצאים בסכנה מיידית), בגלל המצב הקשה בארץ. המדינה המפורסמת ביותר שזה קרה בה הייתה מרוקו. 
	
	בהרבה מקרים ישאלו אותכם בבגרות לגבי איך כל דבר השפיע במדינה שלמדנו עליה. בסיכום של שרית לוב היא המדינה עליה מפרטים. מומלץ לעשות עליה. 
	
	היהודים שהגיעו לארץ הגיעו למחנות פליטים – המעברות. 
	
	\section*{מלחמת יום כיפור}
	רמת הגולן=סוריה. צבאות מצריים וסוריה תקפו את ישראל בהפתעה והצליחו לכבוש שטחים בסיני וברמת הגולן. הצבאות כבשו שטחים בסיני וברמת הגולן. צה''ל גייס את המילואים והצליח להדוף את הסורים מרמת הגולן ולהגיע ליתרון משמעותי על הצבע המצרי בדרום. ''מלחמת יום הכיפורים היא המדינה עם ההצלחה הכי גדולה של מדינת ישראל``. 
	
	בימים הראשונים הופתענו וספגנו אבדות כבדות. בסופה הגענו 101 קילומטר מקהיר ו־50 קילומטר מדמשק. יש ניצחון צבאי. היא נתפשת ככשלון בגלל ההפתעה והמחדל, והמחיר שהמדינה שלמה במלחמה הזו. היו לכך הרבה השפעות בהבט הפנימי. במדינות ערב המלחמה מוצגת כנצחון – בקהיר את מוזיאון שמציג את ''מלחמת אוקטובר`` כנצחון. יש שם בחוץ טנק ישראלי. זה מוצג כנצחון מאותה הסיבה שאצלנו זה נתפס ככשלון – הם הצליחו להפתיע. 
	\subsection*{מה זה בכלל ששת הימים}
	מלחמת ששת־הימים היא אחד הגורמים למלחמת יום הכיפורים. לכן נדבר עליה בקצרה. בחודשים לפני ששת הימים היו מספר גורמים שהובילו אליה. מדינות ערב תפסו את ישראל כחלשה באותה העת בגלל חילופי השלטון בארץ. בסוריה הגבול עבר מעל הכנרת, ורמת הגולן הייתה סורית. דוגמה, היה איזור מפורז בין סוריה לישראל, שהיווה שטח חקלאי. מדינת ישראל פירשה את המושג בפירוז כ''אין נשק`` לעומת סוריה שפרשה את זה כ''אין אנשים``. הסורים ניסו להטות את מי הירדן לפני ששת הימים, הפילו מטוסים, היה בלגנים. בצד המצרי אחרי ולמרות הסכמי מבצע סיני, המצרים חסמו את מצרי טירן (סחר לישראל, השבית את נמל אילת), רכזו כוחות בסיני, וכו'. הייתה תקופה שנקראה ''תקופת ההמתנה`` – תקופה שבה מדינת ישראל מגייסת מילואים, ולא עושה שום דבר. היו לחצים של האמריקקים לא לעשות מכת מנע וליזום מלחמה. בגלל הכניעה ללחצים האמריקאים בתקופה זו, לוי אשכול נתפש כדרדלה, למרות שהוא לבסוף קיבל את ההחלטה לצאת מלחמת ששת הימים. 
	
	מדינת ישראל יצאה למלחמה מול שלוש מדינות – מצריים, ירדן וסוריה. בשורה התחתונה, נצחנו תוך שישה ימים, ומדינת ישראל שלישה את שטחה, וכובשת מעבר לקו הירוק (קווי שביתת הנשק של מלחמת העצמאות). 
	
	
	\ndoc
	
\end{document}